\chapter{What all three methods share}
\label{ch:common}

The three acquisition methods run the \emph{same} FPGA program and read the
\emph{same} accumulated buffer. They differ only in how many repetitions the host
takes per call and how it drains the buffer. Everything in this chapter is
therefore common to all three, and the per-method chapters that follow only have
to describe the differences.

\section{The clock everything is counted in}
\label{sec:clock}

Every duration in the program is stored as an integer number of cycles --- the
\co{_treg} suffix on every \co{NVConfiguration} field. The conversion is done by
\co{soccfg.us2cycles()} and the clock is

\begin{equation}
f_\text{tProc} = \SI{307.2}{\mega\hertz}
\qquad\Longrightarrow\qquad
1\ \text{treg} = \frac{1}{\SI{307.2}{\mega\hertz}} = \SI{3.2552}{\nano\second}.
\label{eq:clock}
\end{equation}

That number is not taken from the firmware printout; it is recovered from the
data, by the argument in \S\ref{sec:quantum}. The recovered divisor for the
\SI{213}{\micro\second} runs is $D=65434$, and
$65434/\num{307.2e6} = \SI{213.001}{\micro\second}$ --- and \co{us2cycles(213)}
rounds $213\times307.2 = 65433.6$ to exactly \num{65434}. The two agree to the
last cycle, which is what makes it safe to quote the rest of this chapter in
nanoseconds.

\begin{center}
\small
\begin{tabular}{@{}lrrl@{}}
\toprule
\hd{Field} & \hd{Value} & \hd{treg} & \hd{Exact duration} \\
\midrule
\co{readout_integration_tus} = 120 (current) & \SI{120}{\micro\second} & \num{36864} & \SI{120.0000}{\micro\second} \\
\co{readout_integration_tus} = 213 (08-06)   & \SI{213}{\micro\second} & \num{65434} & \SI{213.0013}{\micro\second} \\
\co{relax_delay_treg} (set in the run cells) & --- & \num{1000} & \SI{3.2552}{\micro\second} \\
\co{relax_delay_tns} (\co{default_config})   & \SI{500000}{\nano\second} & \num{153600} & \SI{500.0000}{\micro\second} \\
\co{synci(100)} in \co{initialize()}         & --- & \num{100} & \SI{325.52}{\nano\second} \\
\bottomrule
\end{tabular}
\end{center}

\co{default_config.relax_delay_tns = 500000} is \emph{overridden} by every
acquisition cell, each of which sets \co{relax_delay_treg = 1000}. The
\SI{500}{\micro\second} value never reaches the FPGA in any of the three methods;
it would double the rep time if it did.

\section{The FPGA program: one rep, step by step}
\label{sec:asm}

\co{MultipointLockinODMR.body()}
(\file{notebook_modules/multipoint_lockin_program.py}) emits the following
assembly \emph{per parked frequency}, i.e.\ per \emph{emission slot}. The last
column is what each step contributes to the tProc time cursor --- which is not
the same as how long the step lasts, and the difference is the whole content of
this table.

\begin{center}
\small
\begin{tabular}{@{}c>{\raggedright\arraybackslash}p{4.6cm}>{\raggedright\arraybackslash}p{5.6cm}>{\raggedright\arraybackslash}p{2.9cm}@{}}
\toprule
\hd{\#} & \hd{ASM step} & \hd{What it does} & \hd{Advances the cursor by} \\
\midrule
1 & \co{mw_frequency_register.set_to(f)} & Retunes the MW generator at runtime.
    This is the trick that lets 2 (or 16) non-uniform frequencies live in one
    program instead of one upload each. & \textbf{0}. Register writes are tProc
    instructions, not scheduled events. \\[2pt]
2 & \co{pulse(ch=mw_channel, t=0)} & Constant-amplitude MW pulse. Its length is
    \co{readout_integration_treg} --- the MW is on for exactly the readout
    window, not longer. & Sets the \emph{generator} timestamp to
    $0+\tau = \SI{120.0000}{\micro\second}$ \\[2pt]
3 & \co{trigger_no_off(adcs, pins=[laser_gate],} \co{width=readout_integration_treg,}
    \co{adc_trig_offset=0, t=0)}
  & Opens the laser gate \emph{and} arms the ADC accumulator, both at $t=0$.
    \textbf{This step is the rep.} & Sets the \emph{readout} timestamp to
    $0+\tau = \SI{120.0000}{\micro\second}$ \\[2pt]
4 & \co{trigger(pins=[laser_gate],} \co{width=relax_delay_treg,}
    \co{t=readout_integration_treg)}
  & Closing trigger: drives the laser gate low at the end of the slot. Pins only
    --- no \co{adcs}. & \textbf{0}. See the box below. \\[2pt]
5 & \co{wait_all()} & Blocks the tProc until every timestamp is reached. &
    \textbf{0} beyond what 2 and 3 already set \\[2pt]
6 & \co{sync_all(relax_delay_treg)} & New cursor $=$ max timestamp $+$
    \num{1000}\,treg. & \SI{3.2552}{\micro\second} \\
\bottomrule
\end{tabular}
\end{center}

\begin{keybox}[Why step 4's \SI{3.2552}{\micro\second} costs nothing, from the source]
\co{trigger()} only writes a timestamp inside \co{if any([adcs, ddr4, mr]):}
(\file{qickdawg/nvpulsing/nvaverageprogram.py}). Step 4 passes \co{pins} and
nothing else, so that branch is skipped entirely: the call emits two \co{seti}
output-set instructions and touches no channel's timestamp. Its \co{width} is a
property of the laser gate waveform, not of the schedule.

\co{trigger_no_off()} is the same function with the closing \co{seti} commented
out --- which is what the name means. It \emph{does} set the readout timestamp,
to \co{t_start + ro_length} converted into tProc cycles, and that is step 3's
entry above.
\end{keybox}

% \texorpdfstring because a bare \SI{...}{\percent} expands to a literal `%' in the
% PDF bookmark, which comments out the rest of the .out line and breaks the build.
\subsection{\texorpdfstring{The per-slot budget, and the \SI{1.6}{\percent} that is left over}{The per-slot budget, and the 1.6 percent that is left over}}

Summing the last column: one slot advances the cursor by
$\tau + \SI{3.2552}{\micro\second}$, so the programmed rep time is

\begin{equation}
t_\text{rep}^\text{programmed} \;=\; n_\text{slots}\times\left(n_\text{readouts}\,\tau + \SI{3.2552}{\micro\second}\right).
\label{eq:trepprog}
\end{equation}

At the two operating points that have been measured:

\begin{center}
\small
\begin{tabular}{@{}lrrrr@{}}
\toprule
& \hd{$n_\text{slots}\tau$} & \hd{Eq.~\ref{eq:trepprog}} & \hd{measured cadence} & \hd{Eq.~\ref{eq:trepprog}\,$-$\,measured} \\
\midrule
$\tau=\SI{213}{\micro\second}$, 2 slots & \SI{426.0}{\micro\second} & \SI{432.5}{\micro\second} & \SI{423.8}{}--\SI{425.4}{\micro\second} & $+7$ to \SI{+9}{\micro\second} (1.6--2.0\%) \\
$\tau=\SI{120}{\micro\second}$, 2 slots & \SI{240.0}{\micro\second} & \SI{246.5}{\micro\second} & \SI{240.0}{\micro\second}$^\dagger$ & \SI{+6.5}{\micro\second} (2.7\%) \\
\bottomrule
\end{tabular}
\end{center}

\noindent{\footnotesize $^\dagger$recovered exactly from the value quantisation
(\S\ref{sec:quantum}); the \SI{213}{\micro\second} figures are wall-clock cadences
across four 2026-08-06 burst runs.}

The measurement sits \emph{below} even $n_\text{slots}\tau$ at
$\tau=\SI{213}{\micro\second}$, so the two \co{sync_all} terms are not merely
small --- they are not visible at all. The model used everywhere in this document
therefore drops them; \S\ref{sec:trep} states it and \co{time_per_rep()}
implements it.

\subsection{How many slots there are}

\begin{center}
\small
\begin{tabular}{@{}llcc@{}}
\toprule
\hd{\co{multipoint_pair_order}} & \hd{\co{multipoint_skip_reference}} & \hd{Slots/rep} & \hd{Readouts/rep} \\
\midrule
\co{"forward"} & \co{True}  \emph{(current default)} & 2 & \textbf{2} \\
\co{"forward"} & \co{False} & 2 & 4 \\
\co{"abba"}    & \co{True}  & 4 & \textbf{4} \\
\co{"abba"}    & \co{False} & 4 & 8 \\
\bottomrule
\end{tabular}
\end{center}

\co{multipoint_skip_reference = True} drops the per-slot ``MW-off'' reference
readout. It is on because the MW cannot actually be gated off on this setup, so
the reference was a contaminated copy of the signal rather than a clean
normaliser --- it halved the FPGA time and removed a systematic. With it on, $z$
is normalised against a stored constant (\co{ref_norm_counts} from the
calibration JSON) instead of a per-shot reference.

\subsection{The one timing formula everything else derives from}
\label{sec:trep}

\begin{equation}
\boxed{\;t_\text{rep} \;=\; n_\text{slots}\times n_\text{readouts per slot}\times \tau\;}
\label{eq:trep}
\end{equation}

Worked through for every configuration this project runs, at
$\tau=\SI{120}{\micro\second}$:

\begin{center}
\small
\begin{tabular}{@{}lccrrr@{}}
\toprule
\hd{Configuration} & \hd{slots} & \hd{reads/slot} & \hd{$t_\text{rep}$} & \hd{cadence} & \hd{reads/shot} \\
\midrule
two-point, forward, skip ref \emph{(default)} & 2  & 1 & \SI{240}{\micro\second}  & \SI{4167}{\hertz} & 2 \\
two-point, forward, with ref                  & 2  & 2 & \SI{480}{\micro\second}  & \SI{2083}{\hertz} & 4 \\
two-point, ABBA, skip ref                     & 4  & 1 & \SI{480}{\micro\second}  & \SI{2083}{\hertz} & 4 \\
16-point, forward, skip ref                   & 16 & 1 & \SI{1.92}{\milli\second} & \SI{521}{\hertz}  & 16 \\
16-point, forward, with ref \emph{(Step 4S)}  & 16 & 2 & \SI{3.84}{\milli\second} & \SI{260}{\hertz}  & 32 \\
\bottomrule
\end{tabular}
\end{center}

The last column is not decoration: it is what the accumulated buffer budget is
spent on, and \S\ref{sec:buffer} and the streaming chapters both come back to it.

\begin{keybox}[The relax windows do not appear in Eq.~\ref{eq:trep}, and that is measured]
An earlier model added the relax windows, giving
$n_\text{slots}\times(\tau + 2\times\SI{3.2552}{\micro\second})$. Reading the
source narrows that to \emph{one} \co{sync_all} per slot rather than two (the
closing \co{trigger} sets no timestamp --- see the box in \S\ref{sec:asm}), and
the measurement removes even that one: across four 2026-08-06 burst runs the
cadence is \SI{423.8}{}--\SI{425.4}{\micro\second} at
$\tau=\SI{213}{\micro\second}$, against \SI{426.0}{\micro\second} for
$2\times\tau$ and \SI{432.5}{\micro\second} for the sync-inclusive sum. The
measurement is \emph{below} $2\tau$, so the sync term is not resolvable.
\co{time_per_rep()} implements Eq.~\ref{eq:trep}; \co{include_relax=True}
restores the old sum as a conservative upper bound if you want headroom rather
than an estimate.
\end{keybox}

\section{What the accumulated buffer holds}
\label{sec:buffer}

The ADC does not hand back samples; it hands back a \emph{sum}. For each readout
trigger, the firmware accumulates \co{readout_integration_treg} raw ADC samples
into one 64-bit integer entry, and writes one such entry per readout per rep. So
after a run of \co{reps} repetitions with $R$ readouts per rep, the buffer holds
$\text{reps}\times R$ integers, laid out as
\co{(reps, nsweep, R)} --- with \co{nsweep = 1} here.

\co{analyze_results()} then does two divisions:

\begin{equation}
\text{stored value} \;=\; \frac{\text{integer accumulator sum}}
                               {\underbrace{\texttt{readout\_integration\_treg}}_{\text{always}}
                                \times
                                \underbrace{\texttt{reps}}_{\text{averaged mode only}}}
\label{eq:divisions}
\end{equation}

Two consequences, both used heavily later:

\begin{enumerate}
  \item \textbf{The stored counts are $\tau$-independent in the mean} (they are a
  mean ADC level, not a total), which is why the counts do not grow when $\tau$
  is raised. Their \emph{variance} does fall as $1/\tau$, which is the sensitivity
  benefit of a longer window.
  \item \textbf{The exact FPGA readout time is recoverable from the recorded
  numbers.} This is the measurement that settled the timing question; it gets its
  own section.
\end{enumerate}

\section{Recovering the exact readout time from the data}
\label{sec:quantum}

Because the numerator of Eq.~\ref{eq:divisions} is an integer, every stored count
is an integer multiple of
\[
\frac{1}{D},\qquad D \equiv \texttt{treg}\times\texttt{reps}.
\]
Searching for the smallest divisor $D$ that makes \emph{all} of a run's counts
integral therefore recovers $D$ exactly, with no reliance on what the notebook
claims it configured. The readout time follows:

\begin{equation}
\boxed{\;t_\text{ADC per call} \;=\; n_\text{slots}\times \frac{D}{f_\text{clk}},
\qquad f_\text{clk} = \SI{307.2}{\mega\hertz}\;}
\label{eq:recover}
\end{equation}

The clock is pinned independently by two runs:
$\SI{213}{\micro\second}\leftrightarrow\texttt{treg}=65434$ and
$\SI{120}{\micro\second}\leftrightarrow\texttt{treg}=36864$, both giving
\SI{307.2}{\mega\hertz}.

\paragraph{Only the product is identifiable.} $D=36864$ is
$(\texttt{treg}=36864)\times(\texttt{reps}=1)$ and
$(\texttt{treg}=768)\times(\texttt{reps}=48)$ equally well. That is sufficient,
because Eq.~\ref{eq:recover} depends only on $D$. In \emph{burst} mode there is
no rep averaging, so $D$ \emph{is} \co{treg} and the readout window follows
directly --- which is how the 2026-08-14 burst run was confirmed to be at
$\tau=\SI{120.0}{\micro\second}$, exactly as configured.

Implemented as \co{twopoint_postprocess.recover_readout_quantum()} and
\co{readout_seconds()}.

\begin{keybox}[Why this matters practically]
It replaced a 12\% error. The burst cadence had been estimated as
\co{acq_seconds / n_samples} over the real batches, which reads
\SI{269}{\micro\second} against a true \SI{240}{\micro\second} because it
includes the host share of each acquire. That 12\% stretch went straight into the
time axis, the quoted rate, and the frequency scale of every spectrum.
\end{keybox}

\section{The host RPC chain}
\label{sec:rpc}

\co{NVAveragerProgram.acquire()}
(\file{qickdawg/nvpulsing/nvaverageprogram.py}, lines 163--304) issues this
sequence of Pyro remote calls to the board on \textbf{every} call:

\begin{center}
\small
\begin{tabular}{@{}c>{\raggedright\arraybackslash}p{5.4cm}>{\raggedright\arraybackslash}p{7.2cm}@{}}
\toprule
\hd{Stage} & \hd{Call} & \hd{Notes} \\
\midrule
1 & \co{config_all(soc, load_envelopes, reset, load_mem)} &
   Uploads the program. \co{load_mem=True}, so the binary and data memory are
   re-sent on every call. Until 2026-08-12 this raised \co{TypeError} on every
   call (wrong keyword) and fell through to defaults. \\
1a & $\rightarrow$ \co{soc.start_src("internal")} & inside \co{config_all} \\
1b & $\rightarrow$ \co{soc.stop_tproc(lazy=not reset)} &
   \textbf{documented no-op on tProc v1} unless \co{reset=True} \\
1c & $\rightarrow$ \co{load_envelopes}, \co{load_bin_program}, \co{load_mem} &
   re-uploaded each call \\
2 & \co{soc.start_src(start_src)} & again, from \co{acquire} \\
3 & \co{config_bufs(soc, enable_avg=True, enable_buf=False)} & arms the accumulator \\
4 & \co{soc.reload_mem()} & no-op unless tProc v2 \\
5 & \co{soc.start_readout(total_count, ...)} & starts the board-side worker \emph{and} the tProc \\
6 & \co{soc.poll_data()} $\times N$ & \textbf{blocks here while the FPGA runs} \\
7 & Pyro deserialise $\rightarrow$ \co{analyze_results()} reshape & host-side \\
\bottomrule
\end{tabular}
\end{center}

\begin{keybox}[The critical structural fact: stage 6 is where the FPGA time goes]
The host does not wait for the FPGA \emph{after} finishing its own work --- it
waits \emph{inside} stage 6, which is part of the chain. So the FPGA time is
\textbf{concurrent with} the host chain, not additive to it, up to the point where
the FPGA run outlasts the rest of the chain.

That single fact is the whole content of Chapter~\ref{ch:timing}, and it is what
the 2026-08-06 analysis got wrong.
\end{keybox}

\begin{openbox}[OPEN --- the per-stage split has never been measured]
The stage list above is read off the source. The \emph{time} each stage takes can
only be measured on the rig, by wrapping the soc proxy so every remote call is
timed. \file{scripts/profile_twopoint_acquire.py} does exactly that and has never
been run, because the RFSoC is not reachable from the analysis machine. Until it
is, the chain is known only in total (10--11.4\,ms) and not stage by stage.

This is the single largest gap in this document. Run:
\begin{center}\co{python scripts/profile_twopoint_acquire.py --ip 192.168.0.103}\end{center}
\end{openbox}

\section{Python outside the call}

Measured directly as (loop period $-$ \co{acq_seconds}), per batch:

\begin{center}
\small
\begin{tabular}{@{}lrr@{}}
\toprule
\hd{Run} & \hd{Python outside \co{acquire()}} & \hd{Share of period} \\
\midrule
2026-08-14 $\tau$=120, 10 reps & \SI{0.067}{\milli\second} & 0.58\% \\
2026-08-14 $\tau$=120, 23 reps & \SI{0.056}{\milli\second} & 0.49\% \\
2026-08-06 $\tau$=213, 23 reps & \SI{0.058}{\milli\second} & 0.51\% \\
\bottomrule
\end{tabular}
\end{center}

Row-dict construction, the conversion arithmetic and the history append together
cost well under a tenth of a millisecond. \textbf{The Python side is not a
bottleneck in any mode} and the earlier vectorisation of the per-sample
conversion did its job. It is listed here so that it can be ruled out rather than
suspected.

The exception is the CSV write, which is not per-sample and is covered per mode.
